The definition of the magnitude is:
\[
m_I = -2.5\lg F_I + const
\]
where $F_I$ is the flux received from the source. Using the data above, we can
obtain the constant:
\begin{align*}
0.0 &= -2.5\lg(0.83 \times 10^{-11}) + const\\
const &= -27.7
\end{align*}
Thus,
\begin{align*}
m_I &= -2.5\lg F_I - 27.7\\
F_I &= 10^{\frac{m_I + 27.7}{-2.5}}
  = 1.3 \times 10^{-20}\,\mathrm{W\,m^{-2}\,nm^{-1}}
\end{align*}
\hfill(4)

For our star, at an effective wavelength $\lambda_0 = 800$\,nm, using this
flux, the number of photons detected per unit wavelength per unit area is the
flux divided by the energy of a photon with the effective wavelength:
\[
N_I = \frac{1.3 \times 10^{-20}}{hc/\lambda_0}
 = 5.3 \times 10^{-2}\ \mathrm{photons\,s^{-1}\,m^{-2}\,nm^{-1}}
\]
\hfill(3)

Thus the total number of photons detected from the star per second by the 8\,m
Gemini telescope over the $I$ band is
\begin{align*}
N_I(total) &= (\text{tel.\ collecting area}) \times QE
  \times Bandwidth \times N_I\\
&= (\pi \times 4^2) \times 0.4 \times 24 \times N_I\\
&= 26\ \mathrm{photons/s} \approx 30\ \mathrm{photons/s}
\end{align*}
\hfill(3)
